\documentclass[12pt,letterpaper]{article}
\usepackage{graphicx,textcomp}
\usepackage{natbib}
\usepackage{setspace}
\usepackage{fullpage}
\usepackage{color}
\usepackage[reqno]{amsmath}
\usepackage{amsthm}
\usepackage{fancyvrb}
\usepackage{amssymb,enumerate}
\usepackage[all]{xy}
\usepackage{endnotes}
\usepackage{lscape}
\newtheorem{com}{Comment}
\usepackage{float}
\usepackage{hyperref}
\newtheorem{lem} {Lemma}
\newtheorem{prop}{Proposition}
\newtheorem{thm}{Theorem}
\newtheorem{defn}{Definition}
\newtheorem{cor}{Corollary}
\newtheorem{obs}{Observation}
\usepackage[compact]{titlesec}
\usepackage{dcolumn}
\usepackage{tikz}
\usetikzlibrary{arrows}
\usepackage{multirow}
\usepackage{xcolor}
\newcolumntype{.}{D{.}{.}{-1}}
\newcolumntype{d}[1]{D{.}{.}{#1}}
\definecolor{light-gray}{gray}{0.65}
\usepackage{url}
\usepackage{listings}
\usepackage{color}

\definecolor{codegreen}{rgb}{0,0.6,0}
\definecolor{codegray}{rgb}{0.5,0.5,0.5}
\definecolor{codepurple}{rgb}{0.58,0,0.82}
\definecolor{backcolour}{rgb}{0.95,0.95,0.92}

\lstdefinestyle{mystyle}{
	backgroundcolor=\color{backcolour},   
	commentstyle=\color{codegreen},
	keywordstyle=\color{magenta},
	numberstyle=\tiny\color{codegray},
	stringstyle=\color{codepurple},
	basicstyle=\footnotesize,
	breakatwhitespace=false,         
	breaklines=true,                 
	captionpos=b,                    
	keepspaces=true,                 
	numbers=left,                    
	numbersep=5pt,                  
	showspaces=false,                
	showstringspaces=false,
	showtabs=false,                  
	tabsize=2
}
\lstset{style=mystyle}
\newcommand{\Sref}[1]{Section~\ref{#1}}
\newtheorem{hyp}{Hypothesis}

\title{Problem Set 3}
\date{Due: November 13, 2025}
\author{Applied Stats/Quant Methods 1}


\begin{document}
	\maketitle
	\section*{Instructions}
	\begin{itemize}
		\item Please show your work! You may lose points by simply writing in the answer. If the problem requires you to execute commands in \texttt{R}, please include the code you used to get your answers. Please also include the \texttt{.R} file that contains your code. If you are not sure if work needs to be shown for a particular problem, please ask.
	\item Your homework should be submitted electronically on GitHub.
	\item This problem set is due before 23:59 on Thursday November 13, 2025. No late assignments will be accepted.

	\end{itemize}

		\vspace{.25cm}
	
\noindent In this problem set, you will run several regressions and create an add variable plot (see the lecture slides) in \texttt{R} using the \texttt{incumbents\_subset.csv} dataset. Include all of your code.

	\vspace{.5cm}
\section*{Question 1}
\vspace{.25cm}
\noindent We are interested in knowing how the difference in campaign spending between incumbent and challenger affects the incumbent's vote share. 
	\begin{enumerate}
		\item Run a regression where the outcome variable is \texttt{voteshare} and the explanatory variable is \texttt{difflog}.	\vspace{2cm}
		
		\noindent Checking structure of data. Regressing the outcome variable on the explanatory variable and showing the output: 
		
\begin{BVerbatim}
	head(inc.sub)
	str(inc.sub)
	regression1 <- lm(voteshare ~ difflog, data = inc.sub)
	summary(regression1)
\end{BVerbatim}
		
		\item Make a scatterplot of the two variables and add the regression line. 	\vspace{2cm}
		
		Creating a gglpot and adding a regression line. The regression line indicates a positive correlation between campaign spending and incumbent vote share.  
		
		\begin{BVerbatim}
ggplot(data = inc.sub, aes(x = difflog, y = voteshare)) + 
geom_point(size = 2, shape = 22) + 
geom_smooth(method = lm, color = "blue") + 
labs(title = "Difference Campaign spending vs. incumbents vote share", 
x = "Campaign spending", 
y = "Incumbents vote share") + 
theme_bw()
		\end{BVerbatim}
		
		Here we see the scatterplot showing the relationship between campaign spending and incumbent vote share: 
		
		\includegraphics[width=0.5\textwidth]{/Users/mairikachur/Documents/PhD/03. Classes/Quantitative Methods I/PS03/Rplot1.png}
		
		\item Save the residuals of the model in a separate object.	\vspace{1cm}
		
		Saving the residuals of the model in the separate object "residuals1" and having a look at the output. 
		\begin{BVerbatim}
residuals1 <- regression1$residuals
residuals1
		\end{BVerbatim}
		
		\item Write the prediction equation.
		To find the prediction equation, we need to find the relevant coefficients: 
		
		\begin{BVerbatim}
summary(regression1)
		\end{BVerbatim}

This summary gives us the following coefficients: 

		\begin{BVerbatim}
Coefficients:
Estimate Std. Error t value Pr(>|t|)    
(Intercept) 0.579031   0.002251  257.19   <2e-16 ***
difflog     0.041666   0.000968   43.04   <2e-16 ***
		\end{BVerbatim}

With these coefficients, we can now write the prediction equation:  $\gamma$ = $\alpha$ + $\beta$*x.
In other terms: Incumbent vote share = intercept + slope * difflog (or x). 
In other terms: 

\begin{BVerbatim}
Incumbent vote share = 0.57903 + 0.04167*difflog
\end{BVerbatim}

On average, for every one-unit increase of the explanatory variable "difflog",  i.e. campagin spending, the output variable "Incumbent vote share" increases by 0.04167. 

	\end{enumerate}
	
\newpage

\section*{Question 2}
\noindent We are interested in knowing how the difference between incumbent and challenger's spending and the vote share of the presidential candidate of the incumbent's party are related.	\vspace{.25cm}
	\begin{enumerate}
		\item Run a regression where the outcome variable is \texttt{presvote} and the explanatory variable is \texttt{difflog}.	\vspace{1cm}
		
Regressing the outcome variable on the explanatory variable:

\begin{BVerbatim}
regression2 <- lm(presvote ~ difflog, data = inc.sub)
\end{BVerbatim}
		
		\item Make a scatterplot of the two variables and add the regression line. 	\vspace{1cm}

Making a scatterplot and adding the regression line. The regression line shows a positive correlation between campaign spending and the voteshare of the presidential candidate, although the correlation is weaker than in the previous question. 



\begin{BVerbatim}
	
ggplot(data = inc.sub, aes(x = difflog, y = presvote)) + 
geom_point(size = 2, shape = 23) + 
geom_smooth(method = lm, color = "orange") + 
labs(title = "Difference in campaign spending and
presidential candidate vote share", 
x = "Campaign spending", 
y = "Voteshare presidential candidate") + 
theme_bw()
	
\end{BVerbatim}
	
		\includegraphics[width=0.5\textwidth]{/Users/mairikachur/Documents/PhD/03. Classes/Quantitative Methods I/PS03/Rplot2.png}
	
		\item Save the residuals of the model in a separate object.	\vspace{1cm}

Saving the residuals of regression2 in the separate object "residuals2" and having a look at the output. 
		
	\begin{BVerbatim}
residuals2 <- regression2$residuals
residuals2
	\end{BVerbatim}
		
		
		\item Write the prediction equation.

To find the prediction equation, we need to find the relevant coefficients: 

\begin{BVerbatim}
	summary(regression2)
\end{BVerbatim}

This summary gives us the following coefficients: 

\begin{BVerbatim}
	Coefficients:
	Estimate Std. Error t value Pr(>|t|)    
	(Intercept) 0.507583   0.003161  160.60   <2e-16 ***
	difflog     0.023837   0.001359   17.54   <2e-16 ***
\end{BVerbatim}

With these coefficients, we can now write the prediction equation:  $\gamma$ = $\alpha$ + $\beta$*x.
\newline In other terms: Voteshare presidential candidate = intercept + slope * difflog (or x). 
In other terms: 

\begin{BVerbatim}

Incumbent vote share = 0.507583 + 0.023837*difflog

\end{BVerbatim}
		
In other words, on average, for every one-unit increase of the explanatory variable "difflog", i.e. campagin spending, the output variable "voteshare presidential candidate" increases by 0.024. 	
	
	\end{enumerate}
	
	\newpage	
\section*{Question 3}

\noindent We are interested in knowing how the vote share of the presidential candidate of the incumbent's party is associated with the incumbent's electoral success.
	\vspace{.25cm}
	\begin{enumerate}
		\item Run a regression where the outcome variable is \texttt{voteshare} and the explanatory variable is \texttt{presvote}.
		
		Regressing voteshare, the vote share of the presidential candidate of the incumbent's party, on presvote, the incumbents electoral success. 
		
		\begin{BVerbatim}
regression3 <- lm(voteshare ~ presvote, data = inc.sub)
		\end{BVerbatim}
		
		
		
			\vspace{1cm}
		\item Make a scatterplot of the two variables and add the regression line. 

 Making a scatterplot and adding the regression line. The regression line shows a positive correlation between the voteshare of the presidential candidate and the voteshare of the incumbent president's party.


\begin{BVerbatim}
ggplot(data = inc.sub, aes(x = presvote, y = voteshare)) + 
geom_point(size = 2, shape = 2) + 
geom_smooth(method = lm, color = "pink") + 
labs(title = "Difference voteshare incumbent presidents party 
vs voteshare presidential candidate", 
x = "Voteshare presidential candidate", 
y = "Voteshare incumbent presidents' party") + 
theme_bw()
\end{BVerbatim}
		
			\includegraphics[width=0.5\textwidth]{/Users/mairikachur/Documents/PhD/03. Classes/Quantitative Methods I/PS03/Rplot3.png}

			\vspace{5cm}
		\item Write the prediction equation.
		
To make the prediction equation, we need to find the relevant coefficients. To do this: 

\begin{BVerbatim}
regression3
\end{BVerbatim}

We get the following coefficients: 

\begin{BVerbatim}
Coefficients:
(Intercept)     presvote  
0.4413       0.3880 
\end{BVerbatim}

We can now write the prediction equation $\gamma$ = $\alpha$ + $\beta$*x.
\newline Voteshare incumbent presidents' party = Intercept + slope * voteshare presidential candidate
\newline Voteshare incumbent presidents' party = 0.4413 + 0.3880*Voteshare presidential candidate
		
On average, for every one-unit increase of the explanatory variable "presvote", i.e.  the incumbents electoral success, the output variable "voteshare", i.e. the vote share of the presidential candidate of the incumbent's party, increases by 0.3880. 

	\end{enumerate}
	

\newpage	
\section*{Question 4}
\noindent The residuals from part (a) tell us how much of the variation in \texttt{voteshare} is $not$ explained by the difference in spending between incumbent and challenger. The residuals in part (b) tell us how much of the variation in \texttt{presvote} is $not$ explained by the difference in spending between incumbent and challenger in the district.
	\begin{enumerate}
		\item Run a regression where the outcome variable is the residuals from Question 1 and the explanatory variable is the residuals from Question 2.	\vspace{1cm}

To do this, we need to create a new dataframe in which we save the residuals from questions 1 and 2: 

\begin{BVerbatim}
	dataframe <- data.frame(residuals_1 = residuals1, 
	residuals_2 = residuals2)
\end{BVerbatim}
		
		\item Make a scatterplot of the two residuals and add the regression line. 	

Making a scatterplot and adding the regression line. The regression line shows a positive correlation between residuals2 and residuals1.  

\begin{BVerbatim}
	regression4 <- lm(residuals_1 ~ residuals_2, data = dataframe)
	ggplot(data = dataframe, aes(x = residuals_2, y = residuals_1)) + 
	geom_point(size = 2, shape = 2) + 
	geom_smooth(method = lm, color = "green") + 
	labs(title = "Residuals2 regressed on residuals1", 
	x = "residuals 2", 
	y = "residuals 1") + 
	theme_bw()
\end{BVerbatim}

			\includegraphics[width=0.5\textwidth]{/Users/mairikachur/Documents/PhD/03. Classes/Quantitative Methods I/PS03/Rplot4.png}

		\item Write the prediction equation.
To make the prediction equation, we need to find the relevant coefficients. To do this: 

\begin{BVerbatim}
	regression4
\end{BVerbatim}

We get the following coefficients: 

\begin{BVerbatim}
Coefficients:
(Intercept)         	residuals_2  
-0.000000000000000001942   0.256877012700097440145  
\end{BVerbatim}

We can now write the prediction equation $\gamma$ = $\alpha$ + $\beta$*x.
\newline Residuals2 = Intercept + slope * residuals1
\newline Residuals2	 = 0 + 0.256 * residuals1


On average, for every one-unit increase of "residuals1" the output variable "residuals2" increases by 0.3880. 


		
	\end{enumerate}
	
	\newpage	

\section*{Question 5}
\noindent What if the incumbent's vote share is affected by both the president's popularity and the difference in spending between incumbent and challenger? 
	\begin{enumerate}
		\item Run a regression where the outcome variable is the incumbent's \texttt{voteshare} and the explanatory variables are \texttt{difflog} and \texttt{presvote}.
		
\begin{BVerbatim}
	regression5 <- lm(voteshare ~ difflog + presvote, data = inc.sub)
	summary(regression5)
\end{BVerbatim}

This gives us the following coefficients: 

\begin{BVerbatim}
	Coefficients:
	Estimate Std. Error t value            Pr(>|t|)    
	(Intercept) 0.4486442  0.0063297   70.88 <0.0000000000000002 ***
	difflog     0.0355431  0.0009455   37.59 <0.0000000000000002 ***
	presvote    0.2568770  0.0117637   21.84 <0.0000000000000002 ***
\end{BVerbatim}


		\item Write the prediction equation.	\vspace{.5 cm}
		
We can now write the prediction equation $\gamma$ = $\alpha$ + $\beta$*x.
\newline Voteshare = Intercept + $\beta_1$ * difflog + $\beta_2$ * presvote
\newline Voteshare	 = 0.448 + 0.035 * difflog + 0.256 * presvote



		\item What is it in this output that is identical to the output in Question 4? Why do you think this is the case?
	\end{enumerate}

The coefficients of "presvote" in regressions 4 and 5 are the same because both regressions answer the same question in two diffrerent ways: how much the presidents' voteshare influences voteshare for incumbent president if account for campaign spending.


\end{document}
